%! Author = jacopo
%! Date = 2/25/26

% Packages

% Document
\begin{abstract}
    Labeled information has shown itself to be effective in machine learning.
    Unfortunately the supervised learning paradigm doesn't come without its limitations as it is very time and resource consuming.
    There are also situations where labeling is not only hard but too unreliable to produce satisfying results.
    Adopting unsupervised learning techniques such as \textit{contrastive learning}, we propose a new \textit{Foundation Model}
    in the field of electroencephalography designed to work in the domain of Affective Computing, NEEGAVI.
    The approach is based on the assumption that different modalities encode complementary aspects of the same underlying phenomena.
    Beyond traditional EEG measurements, we therefore leverage additional modalities to enrich the learned representations.
    Datasets in the field are typically limited in size so we enhance the model by integrating data from the VATE dataset
    through Knowledge Distillation with further information from Video, Audio and Text modalities.
\end{abstract}